\section*{第 \ref{ch:b2:crustal-melting} 章——地壳熔融：混合岩与花岗岩}

\begin{solution}{exo:b2:crustal-melting:1}
(a) 每 \qty{100}{\gram}：$n(\ce{Al2O3}) = 17.03/101.96 = \qty{0.16703}{mol}$，
$n(\ce{CaO}) = 0.09183$, $n(\ce{Na2O}) = 0.06783$, $n(\ce{K2O}) = 0.03076$.
A/CNK $= 0.16703/0.19042 = 0.877$；A/NK $= 0.16703/0.09859 = 1.694$。
(b) A/CNK $< 1 <$ A/NK：\omterm{def:b2:crustal-melting:asi}{准铝质}。长石无法容纳的钙进入
硅酸钙类矿物：角闪石（或单斜辉石、榍石）。
\end{solution}

\begin{solution}{exo:b2:crustal-melting:2}
(a) 白云母和石榴子石容纳的铝超出了长石所能容纳的部分：\omterm{def:b2:crustal-melting:asi}{过铝质}；该花岗岩
可以是 \omterm{def:b2:crustal-melting:types}{S 型花岗岩}。(b) 角闪石和榍石容纳的钙超出了长石所能容纳的部分：
\omterm{def:b2:crustal-melting:asi}{准铝质}。(c) 霓石和钠质角闪石容纳钠，却不占用
长石中的铝：\omterm{def:b2:crustal-melting:asi}{过碱性}。
\end{solution}

\begin{solution}{exo:b2:crustal-melting:3}
每千米岩石的重量为 $2800\times9.81\times1000 = \qty{2.747e7}{Pa}$，故
$P = 0.2747\,z$（\unit{kbar}，$z$ 以 \unit{km} 计）。固相线为 $T = 680 - 7.5(P -
2) = 695 - 2.060\,z$。令 $10 + 20z = 695 - 2.060z$，得 $z = 685/22.060 =
\qty{31.1}{km}$, $T = \qty{631}{\celsius}$, $P = \qty{8.5}{kbar}$.
\end{solution}

\begin{solution}{exo:b2:crustal-melting:4}
(a) 白色层为\omterm{def:b2:crustal-melting:migmatite}{浅色体}，黑色细带为\omterm{def:b2:crustal-melting:migmatite}{暗色体}，普通片麻岩为
中色体。(b) 边界清楚、界线分明的\omterm{def:b2:crustal-melting:migmatite}{浅色体}、\omterm{def:b2:crustal-melting:migmatite}{暗色体}和中色体：属\omterm{def:b2:crustal-melting:migmatite}{变熔混合岩}，
具条带状构造。(c) 该岩石已变为\omterm{def:b2:crustal-melting:migmatite}{深熔混合岩}（呈条纹状和云雾状）：
熔融进一步发展，暗色矿物也熔化了，熔体仍留在原地，
使层理消失。
\end{solution}

\begin{solution}{exo:b2:crustal-melting:5}
(a) $M = 39.10 + 3\times24.31 + 26.98 + 3\times28.09 + 12\times16.00 + 2\times1.008 =
\qty{417.30}{\gram\per\mole}$; $w = 18.02/417.30 = \qty{4.32}{\percent}$.
(b) $F = 0.25\times0.0432/0.040 = 0.270$，即熔体按质量计为 \qty{27}{\percent}，处于
按体积计 20 至 \qty{30}{\percent} 的范围内，该范围由
相似岩石上的黑云母\omterm{def:b2:crustal-melting:reactions}{脱水熔融}实验得出。
\end{solution}

\begin{solution}{exo:b2:crustal-melting:6}
(a) 钾长石 \qty{278.35}{\gram\per\mole}，钠长石 262.24，钙长石 278.22，白云母
398.33。(b) 物质的量：钾长石 $30/278.35 = \qty{0.10778}{mol}$，钠长石 $0.13347$，
钙长石 $0.01797$，白云母 $0.01255$。钾长石为 $\tfrac12$\ce{K2O.}$\tfrac12$\ce{Al2O3}，
钠长石 $\tfrac12$\ce{Na2O.}$\tfrac12$\ce{Al2O3}，钙长石 \ce{CaO.Al2O3}，白云母
$\tfrac12$\ce{K2O.}$\tfrac32$\ce{Al2O3}（每个化学式单位）。$n(\ce{Al2O3}) = 0.05389 +
0.06673 + 0.01797 + 0.01883 = \qty{0.15742}{mol}$；$n(\ce{K2O}) = 0.05389 + 0.00628 =
0.06017$，$n(\ce{Na2O}) = 0.06673$，$n(\ce{CaO}) = 0.01797$，合计 $0.14487$。
A/CNK $= 0.15742/0.14487 = 1.087$。(c) 若不计白云母，三种长石给出的数值恰为
1（\cref{prop:b2:crustal-melting:feldspar}）；白云母每个化学式单位
比钾长石多一个 \ce{Al2O3}，使该岩石为\omterm{def:b2:crustal-melting:asi}{过铝质}。
\end{solution}

\begin{solution}{exo:b2:crustal-melting:7}
玄武岩：$900 + 2.5\times(1200 - 850) = \qty{1775}{\mega\joule\per\cubic\metre}$。地壳：
$2.5\times(850 - 650) + 0.20\times900 = \qty{680}{\mega\joule\per\cubic\metre}$.
$V_c/V_b = 1775/680 = 2.6$。该估算忽略了热量向上和向两侧传导散失，
以及热量扩散所需的时间；此外还假设玄武岩的全部热量都传给了
已准备好熔融的地壳，因此这是一个上限。
\end{solution}

\begin{solution}{exo:b2:crustal-melting:8}
(a) $n(\ce{Al2O3}) = 14.15/101.96 = 0.13878$，$n(\ce{CaO}) = 0.01753$，$n(\ce{Na2O}) =
0.07809$，$n(\ce{K2O}) = 0.04416$；A/CNK $= 0.13878/0.13978 = 0.993$，A/NK $=
0.13878/0.12225 = 1.135$：勉强属\omterm{def:b2:crustal-melting:asi}{准铝质}。(b) 若 \ce{Al2O3} 为 \qty{14.30}{\percent}，
$0.14025/0.13978 = 1.003$；若为 \qty{14.00}{\percent}，$0.13731/0.13978 = 0.982$。指数
在认证不确定度范围内跨越 1：该黑曜岩位于
\omterm{def:b2:crustal-melting:asi}{准铝质}–\omterm{def:b2:crustal-melting:asi}{过铝质}的分界线上，其类别并不可靠。
\end{solution}

\begin{solution}{exo:b2:crustal-melting:9}
(a) 平均值 $\bar P = \qty{5.0}{kbar}$，$\bar w = \qty{10.045}{\percent}$；
$\sum(P - \bar P)(w - \bar w) = 17.62$，$\sum(P - \bar P)^2 = 14.0$；斜率
$17.62/14.0 = \qty{1.259}{\percent}$ 每千巴，截距 $10.045 - 1.259\times5.0 =
\qty{3.75}{\percent}$。斜率与所引的 1.27 在
\qty{1}{\percent} 以内相符。(b) $3.75 + 1.259\times6 = \qty{11.30}{\percent}$。(c) 熔体量为 $10/0.1130 =
\qty{88.5}{\gram}$。(d) 该直线仅依据 3 至 \qty{8}{kbar} 之间的
\qty{800}{\celsius} 下的数据；无法保证在更高压力下溶解度仍保持线性，
此时熔体与流体趋于相似。
\end{solution}

\begin{solution}{exo:b2:crustal-melting:10}
(a) 该直线在 \qty{50}{\celsius} 内升高 \qty{4}{kbar}：$P = 6 + 0.08\,(T - 750)$；
在 \qty{770}{\celsius} 时，$P = \qty{7.6}{kbar}$。(b) 压力自 7.6 降至 \qty{6}{kbar} 时，
温度沿该直线由 770 降至 \qty{750}{\celsius}；释放的热量
$2.5\times20 = \qty{50}{\mega\joule\per\cubic\metre}$，可熔化 $50/900 = 0.056$，约
按体积计占岩石的 \qty{6}{\percent}。(c) 这正是为喜马拉雅
\omterm{def:b2:crustal-melting:leucogranite}{淡色花岗岩}所提出的窗口，即 6 至 \qty{8}{kbar}、750 至 \qty{770}{\celsius}：仅靠减压
就能产生熔体，但每穿越一次曲线，生成量只有百分之几。
\end{solution}

\begin{solution}{exo:b2:crustal-melting:11}
$\ell^2/\kappa = (2000)^2/10^{-6} = \qty{4e12}{s} = \qty{0.13}{Ma}$，对于\omterm{def:b2:magma-emplacement:bodies}{岩床}；
$(10^4)^2/10^{-6} = \qty{1e14}{s} = \qty{3.2}{Ma}$，对于底侵层。两者均
远短于 1000 万年：单次事件无法使地壳如此长时间保持熔融，这需要
热量反复补给，可以通过多次侵入，或由
增厚地壳的缓慢升温来实现。
\end{solution}

\begin{solution}{exo:b2:crustal-melting:12}
(a) $5.0x + 2.0(1 - x) = 2.75$ 得 $x = 0.75/3.0 = 0.25$：新成体的四分之一为
\omterm{def:b2:crustal-melting:migmatite}{浅色体}。(b) $0.25\times0.20 + 0.75\times4.0 = 0.05 + 3.00 = \qty{3.05}{\percent}$，与
中色体中的值一致：封闭体系的解释成立。(c) 以 \qty{100}{\gram} 母岩计，
形成 \qty{25}{\gram} 熔体和 \qty{75}{\gram} 残留体；若有 $m$ 克
熔体离开，则 $(25 - m)/(100 - m) = 0.10$，故 $25 - m = 10 - 0.1m$，$m = 15/0.9 =
\qty{16.7}{\gram}$：约三分之二的熔体离开了露头。
\end{solution}

\begin{solution}{pb:b2:crustal-melting:1}
\noindent\textbf{第一部分。}

\textbf{1.} $39.10 + 3\times26.98 + 3\times28.09 + 12\times16.00 + 2\times1.008 =
\qty{398.33}{\gram\per\mole}$.

\textbf{2.} $18.02/398.33 = \qty{4.52}{\percent}$.

\textbf{3.} $200\times0.0452 = \qty{9.04}{\gram}$ 水，即每千克片岩含水。

\textbf{4.} $9.97 + (13.82 - 9.97)\times2/3 = \qty{12.54}{\percent}$。最少熔体量为
$9.04/0.1254 = \qty{72.1}{\gram}$ 每千克：$F \ge \qty{7.21}{\percent}$
（见\cref{prop:b2:crustal-melting:water}）。

\textbf{5.} $9.04/0.060 = \qty{151}{\gram}$: $F = \qty{15.1}{\percent}$.

\textbf{6.} 例如：岩石中白云母可能较少；部分白云母可能
已在固相线以下与石英分解为钾长石和铝硅酸盐，
其中的水在熔融开始前以流体形式逸出；熔体中的水可能多于
\qty{6}{\percent}；部分水可能保留在黑云母中。

\noindent\textbf{第二部分。}

\textbf{7.} $680 - 7.5\times(7.8 - 2) = \qty{636.5}{\celsius}$；岩石的峰值温度
比它低 \qty{43.5}{\celsius}，即使有自由水也无法熔融。

\textbf{8.} $605 + (19.2 - 15.3)/(27 - 15.3)\times65 = \qty{626.7}{\celsius}$。
$t = 0$ 时，底部温度为 \qty{595}{\celsius}，低 \qty{32}{\celsius}；
在 \qty{10}{Ma} 时，底部温度为 \qty{656}{\celsius}，已高于它。

\textbf{9.} 湿固相线处的熔融需要自由水，而受热的变质岩含水极少：
自由水很快被第一批熔体耗尽，此后熔体停止增长，
直到云母分解为止。

\textbf{10.} 在 30 至 \qty{40}{Ma} 之间：$30 + (800 - 795)/(845 - 795)\times10 =
\qty{31.0}{Ma}$。

\textbf{11.} 底部在 70 至 \qty{80}{Ma} 之间降至 \qty{12.5}{kbar} 以下：
$70 + (13.7 - 12.5)/(13.7 - 12.4)\times10 = \qty{79.2}{Ma}$，此时温度为
$914 - 14\times0.923 = \qty{901}{\celsius}$，位于角闪石的稳定区间内：底部若为角闪岩，
则将发生熔融。

\textbf{12.} $90 + (11.0 - 10)/(11.0 - 9.6)\times10 = \qty{97.1}{Ma}$，此时温度为 $867 -
53\times0.714 = \qty{829}{\celsius}$。这高于白云母的稳定区间：白云母
已经分解。该温度位于黑云母分解的 800 至 \qty{850}{\celsius} 范围内
（该数值在较低压力下测定）：片岩的另一含水矿物黑云母已接近其
极限。

\noindent\textbf{第三部分。}

\textbf{13.} $n(\ce{Al2O3}) = 15.0/101.96 = \qty{0.14712}{mol}$, $n(\ce{CaO}) =
0.55/56.08 = 0.00981$, $n(\ce{Na2O}) = 3.4/61.98 = 0.05486$, $n(\ce{K2O}) = 4.9/94.20 =
0.05202$.

\textbf{14.} A/CNK $= 0.14712/0.11669 = 1.261$.

\textbf{15.} A/NK $= 0.14712/0.10688 = 1.377$；A/CNK $> 1$：\omterm{def:b2:crustal-melting:asi}{过铝质}。

\textbf{16.} \omterm{def:b2:crustal-melting:asi}{过铝质}，源自片岩：为 \omterm{def:b2:crustal-melting:types}{S 型花岗岩}。其过剩的铝赋存于
白云母和石榴子石（或堇青石）中。

\textbf{17.} $0.14712 - 0.11669 = \qty{0.03043}{mol}$ 的 \ce{Al2O3}；以白云母计，
$0.03043\times398.33 = \qty{12.1}{\gram}$，以 \qty{100}{\gram} 花岗岩计。

\textbf{18.} 该黑曜岩（恰好）为\omterm{def:b2:crustal-melting:asi}{准铝质}，不具有变沉积熔体的过剩铝，
其源岩可能不同。但单一样品说明不了太多问题：
每个花岗质岩系都含有轻微\omterm{def:b2:crustal-melting:asi}{过铝质}的成员，而岩系须依其整体
趋势来判断。

\noindent\textbf{第四部分。}

\textbf{19.} $10^{12}\,\unit{\cubic\metre}\times2650 = \qty{2.65e15}{kg}$.

\textbf{20.} 花岗岩即熔体减去其中的水：$2.65\times10^{15}/0.94 =
\qty{2.82e15}{kg}$ 熔体。

\textbf{21.} 片岩中仅 $0.151\times0.70$ 的质量以熔体形式离开：
$2.82\times10^{15}/(0.151\times0.70) = \qty{2.67e16}{kg}$.

\textbf{22.} $2.67\times10^{16}/2800 = \qty{9.5e12}{\cubic\metre}$，即
\qty{9.5e3}{\cubic\kilo\metre}，接近该\omterm{def:b2:magma-emplacement:bodies}{深成岩体}的十倍。

\textbf{23.} $9.5\times10^3/5 = \qty{1.9e3}{\square\kilo\metre}$，相当于边长约
\qty{44}{km} 的正方形。

\textbf{24.} 取 $F = \qty{7.21}{\percent}$ 时，源区较大 $15.1/7.21 = 2.1$ 倍，
约 \qty{2.0e4}{\cubic\kilo\metre}。形成该\omterm{def:b2:magma-emplacement:bodies}{深成岩体}所熔融的源区地壳体积，
以第~5~题的\omterm{def:b2:mantle-melting:fraction}{熔体分数}计，为
\textbf{\qty[text-series-to-math=true]{9.5e3}{\cubic\kilo\metre}}。
\end{solution}
